% ==========================================================================
%  Chapter 52 — Material Module: Architectural Analysis & Improvement Plan
% ==========================================================================

\chapter{Material Module: Architectural Analysis and Improvement Plan}
\label{ch:material-arch-improvements}

After completing the six convergence phases (Chapters~\ref{ch:convergence-master-plan}--\ref{ch:phase6-mpi-scalability}), we turn to a deep architectural analysis of the \texttt{Material} module.
This chapter documents the six-layer architecture, identifies semantic and structural deficiencies, presents a phased improvement plan, and records the first implemented fix.


%% ------------------------------------------------------------------
\section{Module Architecture Overview}
\label{sec:mat-arch-overview}

The material module is organized in six conceptual layers (L0--L5),
each building on the previous one.

\begin{description}
  \item[L0 --- Constitutive Concepts.]
    C\texttt{++}20 concepts
    (\texttt{KinematicMeasure}, \texttt{TensionalConjugate}, \texttt{ConjugatePair})
    defined in \texttt{ConstitutiveRelation.hh}
    constrain the strain--stress type pairings at compile time.

  \item[L1 --- Constitutive Relation Hierarchy.]
    Three concept tiers:
    \texttt{ConstitutiveRelation} (general),
    \texttt{ElasticConstitutiveRelation} (path-independent),
    \texttt{InelasticConstitutiveRelation} (history-dependent).
    Concrete classes: \texttt{ElasticRelation<Policy>},
    \texttt{ContinuumIsotropicRelation},
    \texttt{TimoshenkoBeamSection3D/2D},
    \texttt{MindlinShellSection},
    \texttt{PlasticityRelation<\ldots>},
    \texttt{HyperelasticRelation},
    \texttt{FiniteStrainDamageRelation}.

  \item[L2 --- Constitutive Site.]
    \texttt{MaterialInstance<Relation, StatePolicy>}
    pairs a relation (shared ownership via \texttt{shared\_ptr}) with a
    committed state (\texttt{CommittedConstitutiveState}).

  \item[L3 --- Integration Algorithm.]
    \texttt{ConstitutiveIntegrator},
    \texttt{PassThroughIntegrator},
    \texttt{ElasticUpdate} ---
    algorithmic wrappers that select the local iteration strategy
    (forward map for elasticity, return-mapping for plasticity).

  \item[L4 --- Type-Erased Handle.]
    \texttt{Material<Policy>} (owning, heap),
    \texttt{MaterialRef<Policy>} (non-owning, SBO 24\,B),
    \texttt{MaterialConstRef<Policy>} (read-only ref) ---
    External Polymorphism pattern via \texttt{if constexpr} dispatch.

  \item[L5 --- Integration-Point Binding.]
    \texttt{MaterialPoint} and \texttt{MaterialSection} bind a type-erased
    material handle to a quadrature point, supplying the kinematic
    field (strain/generalized-strain) and retrieving the conjugate
    response and tangent.
\end{description}


%% ------------------------------------------------------------------
\section{Identified Deficiencies}
\label{sec:mat-deficiencies}

Table~\ref{tab:deficiencies} summarizes the eight deficiencies
found during the architectural review.

\begin{table}[htbp]
\centering
\small
\begin{tabular}{clcc}
\toprule
  \textbf{\#} & \textbf{Deficiency} & \textbf{Severity} & \textbf{Complexity} \\
\midrule
  1 & \texttt{compliance\_matrix\_} stores \emph{stiffness} $\mathbf{C}$, not compliance $\mathbf{S}=\mathbf{C}^{-1}$ & Medium & Low \\
  2 & Mutable cache in \texttt{PlasticityRelation} (thread-safety) & High & Medium \\
  3 & Triple \texttt{if constexpr} dispatch $\times$18 in \texttt{Material.hh} & Medium & Medium \\
  4 & \texttt{StateVariableT $\equiv$ KinematicT} (state = strain, not internal vars) & High & High \\
  5 & \texttt{shared\_ptr<Relation>} even for exclusive ownership & Low & Low \\
  6 & 8\,M heap allocations for large meshes via \texttt{Material<>} & Medium & High \\
  7 & \texttt{Corotational} placeholder in \texttt{KinematicPolicy.hh} & Functional & Medium \\
  8 & \texttt{BeamElement} lacks \texttt{KinematicPolicy} template param & Low & Medium \\
\bottomrule
\end{tabular}
\caption{Material module deficiencies ranked by severity and implementation complexity.}
\label{tab:deficiencies}
\end{table}


%% ------------------------------------------------------------------
\section{Six-Phase Improvement Plan}
\label{sec:improvement-plan}

Based on the deficiency analysis, we define six implementation phases
ordered by risk (ascending) and impact (descending).

\begin{enumerate}
  \item \textbf{Semantic rename:}
        \texttt{compliance\_matrix\_} $\to$ \texttt{stiffness\_matrix\_}
        across the entire elastic hierarchy.
        Purely mechanical, zero behavioral change, immediate clarity gain.

  \item \textbf{Thread safety:}
        Eliminate the \texttt{mutable} cache in \texttt{PlasticityRelation}
        by moving cached stress/tangent into
        \texttt{CommittedConstitutiveState}.

  \item \textbf{Dispatch deduplication:}
        Replace the $3\times18$ \texttt{if constexpr} blocks in
        \texttt{Material.hh} with a CRTP or concept-constrained
        forwarding layer.

  \item \textbf{Kinematic policy on structural elements:}
        Add \texttt{KinematicPolicy} template parameter to
        \texttt{BeamElement} and \texttt{ShellElement}, enabling
        compile-time selection of small-rotation vs.\ corotational
        formulations.

  \item \textbf{Corotational beam formulation:}
        Implement the CR formulation for 3D Timoshenko beams
        (rotation extraction, force transformation, consistent tangent).

  \item \textbf{State separation and arena allocator:}
        Long-term: decouple internal state variables from kinematic
        measures and introduce arena-based allocation for
        \texttt{MaterialInstance} to amortize heap overhead.
\end{enumerate}


%% ------------------------------------------------------------------
\section{Phase 1 Implementation: Semantic Rename}
\label{sec:phase1-rename}

\subsection{Problem Statement}

The protected member \texttt{compliance\_matrix\_} in
\texttt{ElasticRelation<Policy>} stores the \emph{stiffness} tensor
$\mathbf{C}$ (mapping $\boldsymbol{\varepsilon}\mapsto
\boldsymbol{\sigma}=\mathbf{C}\,\boldsymbol{\varepsilon}$), not the
compliance tensor $\mathbf{S}=\mathbf{C}^{-1}$.
Every derived class (\texttt{ContinuumIsotropicRelation},
\texttt{TimoshenkoBeamSection3D/2D}, \texttt{MindlinShellSection}) and
one test file propagated this misnomer.

\subsection{Scope of Changes}

The rename touched \textbf{five files} totaling \textbf{49 occurrences}:

\begin{enumerate}
  \item \textbf{\texttt{ElasticRelation.hh}} (base class):
    \begin{itemize}
      \item Protected member: \texttt{compliance\_matrix\_} $\to$ \texttt{stiffness\_matrix\_}
      \item Public accessor: \texttt{compliance\_matrix()} $\to$ \texttt{stiffness\_matrix()}
      \item All internal uses in \texttt{compute\_response()}, \texttt{tangent()},
            and \texttt{compute\_stress()}
      \item Comment corrections (``tangent (compliance/stiffness)'' $\to$ ``tangent stiffness'')
    \end{itemize}

  \item \textbf{\texttt{IsotropicRelation.hh}}:
    12 assignments in \texttt{update\_elasticity()} plus comment fix.

  \item \textbf{\texttt{TimoshenkoBeamSection.hh}}:
    3D section (8 assignments, \texttt{set\_matrix()}, print accessor)
    and 2D section (4 assignments, print accessor).

  \item \textbf{\texttt{MindlinShellSection.hh}}:
    13 assignments in \texttt{update\_section\_stiffness()},
    \texttt{set\_matrix()}, and print accessor.

  \item \textbf{\texttt{tests/test\_element\_geometry.cpp}}:
    Single call site \texttt{sec.compliance\_matrix()} $\to$
    \texttt{sec.stiffness\_matrix()}.
\end{enumerate}

\subsection{Verification}

After the rename:
\begin{itemize}
  \item \textbf{Build:} 30/30 targets compiled with \texttt{-Werror}, zero warnings.
  \item \textbf{CTest:} 45/47 passed (2 pre-existing failures:
        \texttt{cantilever\_benchmark} and \texttt{snes\_gmsh\_pipeline},
        both due to missing mesh files).
  \item \textbf{Grep:} zero remaining occurrences of
        \texttt{compliance\_matrix} in the source tree.
\end{itemize}


%% ------------------------------------------------------------------
\section{Summary}
\label{sec:mat-improvements-summary}

This chapter established a reference architecture map for the
\texttt{Material} module, catalogued eight deficiencies, and defined
a six-phase improvement roadmap.
Phase~1 (semantic rename, \S\ref{sec:phase1-rename}),
Phase~2 (thread safety, \S\ref{sec:phase2-thread-safety}),
Phase~3 (kinematic dispatch consolidation, \S\ref{sec:phase3-dispatch}),
Phase~4 (kinematic policy on structural elements, \S\ref{sec:phase4-kinematic-policy}), and
Phase~5 (corotational beam formulation, \S\ref{sec:phase5-corotational})
have been implemented and verified.
Phase~6 (state separation and arena allocator) remains for future work
and will follow the same pattern: implement, build with \texttt{-Werror},
run full CTest regression, and document in an appendix to this chapter.


%% ------------------------------------------------------------------
\section{Phase 2 Implementation: Thread-Safe Constitutive Relations}
\label{sec:phase2-thread-safety}

\subsection{Problem Statement}

Three inelastic constitutive relation classes used \texttt{mutable}
member variables to implement a performance cache: after
\texttt{compute\_response($\varepsilon$)} populates the stress and
tangent, a subsequent call to \texttt{tangent($\varepsilon$)} with the
same strain can return the cached tangent without re-running the
integration algorithm.

The affected classes were:
\begin{itemize}
  \item \texttt{PlasticityRelation<Policy, YieldF, Hardening, Flow, YieldFn, ReturnAlg>}
  \item \texttt{KentParkConcrete}
  \item \texttt{MenegottoPintoSteel}
\end{itemize}

Each declared four \texttt{mutable} cache members:
\begin{lstlisting}[language=C++, basicstyle=\ttfamily\small]
mutable ConjugateT  last_stress_{};
mutable TangentT    last_tangent_{};
mutable bool        cache_valid_{false};
mutable StrainVecT  last_strain_cache_{...};
\end{lstlisting}

This pattern causes two problems:
\begin{enumerate}
  \item \textbf{Thread safety:}
        Methods marked \texttt{const} mutate shared state without
        synchronization.  In a parallel assembly loop, two threads
        evaluating the same relation concurrently could produce data
        races (undefined behavior under the C++ memory model).

  \item \textbf{Semantic dishonesty:}
        A \texttt{const} method that writes to \texttt{mutable}
        members violates the principle that \texttt{const} means
        ``logically observable state does not change.''
        The cache \emph{is} observable (it changes the return path
        of \texttt{tangent()}).
\end{enumerate}

\subsection{Why the Cache Was Unnecessary}

All three classes satisfy \texttt{ExternallyStateDrivenConstitutiveRelation}
(Level~3 concept).  When used through \texttt{MaterialInstance<Relation>},
the 2-argument overloads are selected:
\begin{lstlisting}[language=C++, basicstyle=\ttfamily\small]
// MaterialInstance:
if constexpr (ExternallyStateDrivenConstitutiveRelation<Relation>) {
    return relation_->compute_response(k, this->algorithmic_state());
    return relation_->tangent(k, this->algorithmic_state());
}
\end{lstlisting}
These 2-argument overloads are purely functional---they never read or
write the cache.  The 1-argument overloads (which do use the cache)
are legacy fallbacks that are never exercised in production.

\subsection{Scope of Changes}

For each of the three classes:
\begin{enumerate}
  \item Removed the four \texttt{mutable} cache member declarations.
  \item Simplified \texttt{compute\_response(const KinematicT\&)} to
        call the integration kernel and return the result directly,
        without writing to the cache.
  \item Simplified \texttt{tangent(const KinematicT\&)} to always
        recompute (no cache-hit check).
  \item Removed \texttt{cache\_valid\_ = false} from \texttt{update()}.
\end{enumerate}

The 2-argument overloads and the \texttt{commit()} method were
already cache-free and required no changes.

\subsection{Verification}

\begin{itemize}
  \item \textbf{Build:} 8 recompiled targets, zero warnings (\texttt{-Werror}).
  \item \textbf{CTest:} 45/47 passed (same 2 pre-existing failures).
  \item \textbf{Grep:} zero remaining \texttt{mutable} in
        \texttt{PlasticityRelation.hh},
        \texttt{KentParkConcrete.hh}, and
        \texttt{MenegottoPintoSteel.hh}.
  \item Key plasticity tests passed:
        \texttt{plasticity\_composition},
        \texttt{material\_strategy},
        \texttt{phase4\_hysteretic\_cycles},
        \texttt{rc\_fiber\_frame\_nonlinear},
        \texttt{vertical\_slice\_gate},
        \texttt{nonlinear\_validation}.
\end{itemize}


%% ------------------------------------------------------------------
\section{Phase 3 Implementation: Kinematic Dispatch Consolidation}
\label{sec:phase3-dispatch}

\subsection{Problem Statement}

The type-erased model classes in \texttt{Material.hh}
(\texttt{OwningMaterialModel}, \texttt{NonOwningMaterialConstModel},
\texttt{NonOwningMaterialModel}) each override
\texttt{compute\_response()}, \texttt{tangent()}, and
\texttt{commit()} for the full
\texttt{ConstitutiveKinematics<dim>} argument.  Every override
repeated an identical \texttt{if constexpr} block:

\begin{verbatim}
if constexpr (requires { integrator.method(material, kin); }) {
    return integrator.method(material, kin);
} else {
    return integrator.method(
        material,
        continuum::make_kinematic_measure<StateVariableT>(kin));
}
\end{verbatim}

\noindent
This check prefers the continuum-aware integrator overload when
available and falls back to the reduced kinematic measure otherwise.
The pattern appeared \textbf{eight times} across the three model
classes, tripling the maintenance burden and the risk of
inconsistent modifications.

\subsection{Solution: Dispatch Helper Functions}

Three free-function templates were introduced inside
\texttt{namespace impl}, immediately before the model class
definitions:

\begin{description}
  \item[\texttt{dispatch\_response}] Dispatches
    \texttt{integrator.compute\_response(material, kin)}.
  \item[\texttt{dispatch\_tangent}] Dispatches
    \texttt{integrator.tangent(material, kin)}.
  \item[\texttt{dispatch\_commit}] Dispatches
    \texttt{integrator.commit(material, kin)} (\texttt{void} return).
\end{description}

\noindent
Each helper is templated on \texttt{StateVariableT},
\texttt{IntegratorT}, \texttt{MaterialT}, and the spatial dimension
\texttt{Dim}.  The callers pass their integrator and material by
reference (value semantics for the owning model, dereferenced
pointers for the non-owning models):

\begin{verbatim}
// OwningMaterialModel
return dispatch_response<StateVariableT>(integrator_, material_, kin);

// NonOwningMaterialModel
return dispatch_response<StateVariableT>(*integrator_, *material_, kin);
\end{verbatim}

\subsection{Scope of Changes}

\begin{itemize}
  \item \textbf{File modified:} \texttt{src/materials/Material.hh}.
  \item \textbf{Lines removed:} 8 \texttt{if constexpr} blocks
        ($\approx$72 lines of boilerplate).
  \item \textbf{Lines added:} 3 dispatch helpers
        ($\approx$45 lines, single source of truth).
  \item \textbf{Net reduction:} $\approx$27 lines.
  \item No changes to the public API, concept hierarchy, or
        \texttt{ConstitutiveIntegrator} implementations.
\end{itemize}

\subsection{Verification}

\begin{itemize}
  \item \textbf{Build:} 12 recompiled targets, zero warnings
        (\texttt{-Werror}).
  \item \textbf{CTest:} 45/47 passed (same 2 pre-existing failures).
  \item All continuum-path tests passed:
        \texttt{snes\_hyperelastic},
        \texttt{tl\_integration},
        \texttt{updated\_lagrangian},
        \texttt{finite\_strain\_damage},
        \texttt{continuum\_local\_problem}.
  \item All structural-path tests passed:
        \texttt{beam\_element},
        \texttt{beam\_n\_element},
        \texttt{rc\_fiber\_frame\_nonlinear},
        \texttt{nonlinear\_validation}.
\end{itemize}


%% ------------------------------------------------------------------
\section{Phase 4 Implementation: Kinematic Policy on Structural Elements}
\label{sec:phase4-kinematic-policy}

\subsection{Problem Statement}

The \texttt{BeamElement} class hard-coded its kinematic assumptions:
the element frame (rotation matrix $\mathbf{R}$ and length $L$) was
computed once from the reference configuration and never updated,
the global-to-local transformation was a fixed block-diagonal
rotation $\mathbf{T}=\mathrm{blockdiag}(\mathbf{R}^T,
\mathbf{R}^T)$, and no geometric stiffness contribution was
possible.

This prevented compile-time selection of alternative formulations
(e.g., corotational or updated-Lagrangian beams) without duplicating
the entire element class---a gap already solved for continuum elements
by the \texttt{KinematicPolicy} template parameter in
\texttt{ContinuumElement<MaterialPolicy, ndof, KinematicPolicy>}.

\subsection{Design: \texttt{BeamKinematicPolicyConcept}}

A new header \texttt{src/elements/BeamKinematicPolicy.hh} introduces
a minimal concept and two concrete policies:

\begin{description}
  \item[\texttt{beam::BeamKinematicPolicyConcept}.]
    A policy type \texttt{P} must expose two \texttt{static constexpr
    bool} traits:
    \begin{itemize}
      \item \texttt{P::is\_geometrically\_linear} --- controls whether
            the element frame is fixed or updated each iteration.
      \item \texttt{P::needs\_geometric\_stiffness} --- controls whether
            a geometric stiffness $\mathbf{K}_\sigma$ is appended to
            the tangent.
    \end{itemize}

  \item[\texttt{beam::FrameData<Dim>}.]
    A lightweight result struct carrying $\mathbf{R}\in\mathbb{R}^{d
    \times d}$ and $L\in\mathbb{R}$.

  \item[\texttt{beam::SmallRotation}.]
    Default policy; preserves the original behavior:
    \begin{itemize}
      \item \texttt{compute\_frame<Dim>(x$_0$, x$_1$)}: one-time
            computation from reference node coordinates.  2D:
            $\mathbf{e}_1=(x_1-x_0)/L$, $\mathbf{e}_2=$ 90\textdegree{}
            rotation.  3D: Gram--Schmidt with a reference vector.
      \item \texttt{update\_frame}: identity (returns the reference
            frame unchanged).
      \item \texttt{extract\_local\_dofs}: $\mathbf{u}_{\text{loc}} =
            \mathbf{T}\,\mathbf{u}_{\text{glob}}$.
    \end{itemize}

  \item[\texttt{beam::Corotational}.]
    Crisfield-style CR policy (Phase~5, below).
\end{description}

Both policies satisfy a \texttt{static\_assert} against the concept
at the point of definition.

\subsection{Template Signature Change}

The \texttt{BeamElement} class template gains a third parameter with
a default:

\begin{verbatim}
template <typename BeamPolicy,
          std::size_t Dim = BeamPolicy::dim,
          typename KinematicPolicy = beam::SmallRotation,
          typename AsmPolicy = assembly::DirectAssembly>
class BeamElement { ... };
\end{verbatim}

\noindent
Because \texttt{KinematicPolicy} has a default, all existing
instantiations compile unchanged:

\begin{verbatim}
BeamElement<TimoshenkoBeam2D, 2>   // SmallRotation, DirectAssembly
BeamElement<TimoshenkoBeam3D, 3>   // SmallRotation, DirectAssembly
\end{verbatim}

\subsection{Internal Refactoring}

The two formerly private members \texttt{R\_} and \texttt{length\_}
were replaced by two \texttt{beam::FrameData<Dim>} structs:

\begin{description}
  \item[\texttt{frame0\_}:] Reference-configuration frame, computed
    once in \texttt{compute\_frame()} and immutable thereafter.
  \item[\texttt{frame\_n\_}:] Current-configuration frame.  For
    \texttt{SmallRotation} this is always identical to
    \texttt{frame0\_}; for \texttt{Corotational} it is recomputed
    at each Newton iteration via \texttt{update\_current\_frame()}.
\end{description}

\noindent
Key method changes:

\begin{itemize}
  \item \texttt{compute\_frame()} delegates to
    \texttt{KinematicPolicy::compute\-\_frame<Dim>(x$_0$, x$_1$)}
    and stores the result in both \texttt{frame0\_} and
    \texttt{frame\_n\_}.

  \item A new helper \texttt{update\_current\_frame(u\_e)} extracts
    nodal coordinates and displacements, calls
    \texttt{KinematicPolicy::update\_frame<Dim>()}, and stores the
    result in \texttt{frame\_n\_}.

  \item \texttt{compute\_internal\_forces()} now conditionally calls
    \texttt{update\_current\_frame()} via
    \texttt{if constexpr (!KinematicPolicy::is\_geometrically\_linear)},
    followed by
    \texttt{KinematicPolicy::extract\_local\_dofs()} instead of a
    raw $\mathbf{T}\cdot\mathbf{u}$.

  \item \texttt{inject\_tangent\_stiffness()} follows the same guard
    for the frame update and appends the geometric stiffness when
    \texttt{KinematicPolicy::needs\_geometric\_stiffness} is true.

  \item \texttt{commit\_material\_state()} mirrors the same
    conditional frame update and local-DOF extraction.

  \item Shape-function derivatives \texttt{dN1\_ds()}, \texttt{dN2\_ds()}
    use \texttt{frame\_n\_.length} so that \texttt{SmallRotation}
    keeps the reference length and \texttt{Corotational} uses the
    current length.

  \item Accessors \texttt{element\_length()} and
    \texttt{rotation\_matrix()} now read from \texttt{frame\_n\_}.
    New accessors \texttt{reference\_frame()} and
    \texttt{current\_frame()} expose the full \texttt{FrameData}
    bundles.
\end{itemize}

\subsection{Verification}

\begin{itemize}
  \item \textbf{Build:} 44 recompiled targets, zero warnings
        (\texttt{-Werror}, clang++-20, C++23).
  \item \textbf{CTest:} 45/47 passed (same 2 pre-existing mesh-file
        failures).
  \item All beam tests passed:
        \texttt{beam\_element},
        \texttt{beam\_n\_element},
        \texttt{rc\_fiber\_frame\_nonlinear},
        \texttt{nonlinear\_validation}.
\end{itemize}


%% ------------------------------------------------------------------
\section{Phase 5 Implementation: Corotational Beam Formulation}
\label{sec:phase5-corotational}

\subsection{Formulation Overview}

The corotational (CR) approach decomposes element motion into a
rigid-body component and a deformational component.  At each Newton
iteration:

\begin{enumerate}
  \item \textbf{Current chord.}  Compute the deformed node positions
    $\mathbf{x}_i^n = \mathbf{X}_i + \mathbf{u}_i$ and form the
    current chord direction $\mathbf{e}_1^n = (\mathbf{x}_1^n -
    \mathbf{x}_0^n)/l_n$, where $l_n = \|\mathbf{x}_1^n -
    \mathbf{x}_0^n\|$ is the current length.

  \item \textbf{Corotated frame.}  Build the rotation matrix
    $\mathbf{R}^n$ from $\mathbf{e}_1^n$ and a smoothly evolved
    transverse axis.  In 3D, the reference $\mathbf{e}_2^0$ is
    projected onto the plane perpendicular to $\mathbf{e}_1^n$ and
    renormalized.

  \item \textbf{Deformational DOFs.}  Extract local DOFs by
    filtering the rigid rotation of the chord:
    \begin{description}
      \item[2D ($6$ DOFs):] $u_1^{\text{def}} = 0$,
        $v_1^{\text{def}} = 0$,
        $\theta_1^{\text{def}} = \theta_1 - \alpha$ (rigid chord
        rotation $\alpha$),
        $u_2^{\text{def}} = l_n - l_0$,
        $v_2^{\text{def}} = 0$,
        $\theta_2^{\text{def}} = \theta_2 - \alpha$.
      \item[3D ($12$ DOFs):] Translate into the corotated frame;
        subtract the incremental rigid rotation from nodal rotations
        via $\boldsymbol{\theta}_i^{\text{loc}} = \mathbf{R}^n
        \boldsymbol{\theta}_i^{\text{glob}} -
        \boldsymbol{\delta\theta}$, where $\boldsymbol{\delta\theta}$
        is the axial vector of $\mathbf{R}^n(\mathbf{R}^0)^T$.
    \end{description}

  \item \textbf{Local response.}  Evaluate the small-rotation
    Timoshenko formulation with the deformational DOFs and the
    current length $l_n$ (the B matrix shape-function derivatives
    use $l_n$ instead of $l_0$).

  \item \textbf{Global forces.}  Transform back via the corotational
    projector.

  \item \textbf{Geometric stiffness.}  Append
    $\mathbf{K}_\sigma$ from the axial force $N$:
    \[
        (\mathbf{K}_\sigma)_{ij} =
        \frac{N}{l_n}\,\delta_{v_i v_j}\;,
    \]
    acting on the transverse DOF pairs $(v_1, v_2)$ in 2D and
    $(u_{y_1}, u_{y_2})$, $(u_{z_1}, u_{z_2})$ in 3D.
\end{enumerate}

\noindent
References: Crisfield, \emph{Non-Linear Finite Element Analysis of
Solids and Structures}, Vol.~1, \S7.3; Battini \& Pacoste (2002).

\subsection{Implementation in \texttt{beam::Corotational}}

The policy provides four static template functions:

\begin{description}
  \item[\texttt{compute\_frame<Dim>}:]
    Delegates to \texttt{SmallRotation::compute\_frame} for the
    initial (reference) frame computation.

  \item[\texttt{update\_frame<Dim>}:]
    Computes $\mathbf{x}_i^n = \mathbf{X}_i + \mathbf{u}_i$ and
    rebuilds the element chord.  In 3D, the reference $\mathbf{e}_2$
    is projected and renormalized.  Degeneracy is handled by falling
    back to the reference-vector algorithm.

  \item[\texttt{extract\_local\_dofs<Dim, TotalDofs>}:]
    Filters the rigid chord rotation:
    \begin{itemize}
      \item 2D: computes $\alpha = \arctan(\sin\alpha,\cos\alpha)$
            from $\mathbf{e}_1^n\cdot\mathbf{e}_1^0$ and the
            cross-product, then $\theta_i^{\text{def}} =
            \theta_i^{\text{glob}} - \alpha$.
      \item 3D: computes $\mathbf{R}_{\text{rigid}} =
            (\mathbf{R}^n)^T\,\mathbf{R}^0$, extracts the axial
            vector $\boldsymbol{\delta\theta}$ from the
            skew-symmetric part, and subtracts it from each node's
            rotations expressed in the corotated frame.
    \end{itemize}

  \item[\texttt{geometric\_stiffness<Dim, TotalDofs, NumStrains>}:]
    Builds $\mathbf{K}_\sigma$ from the axial force component of the
    section force vector.  2D: $4\times4$ sub-block on
    $(v_1,v_2)$.  3D: two $4\times4$ sub-blocks on $(u_{y_1},
    u_{y_2})$ and $(u_{z_1}, u_{z_2})$.
\end{description}

\subsection{Integration with \texttt{BeamElement}}

The \texttt{BeamElement} methods use \texttt{if constexpr} guards
keyed to the policy traits to enable CR-specific code paths at zero
cost when \texttt{SmallRotation} is selected:

\begin{verbatim}
if constexpr (!KinematicPolicy::is_geometrically_linear) {
    update_current_frame(u_e);  // recompute R_n, l_n
}
auto u_def = KinematicPolicy::template
    extract_local_dofs<Dim, total_dofs>(...);
...
if constexpr (KinematicPolicy::needs_geometric_stiffness) {
    K_local += KinematicPolicy::template
        geometric_stiffness<Dim, total_dofs, num_strains>(...);
}
\end{verbatim}

\noindent
For \texttt{SmallRotation}, the compiler eliminates these branches
entirely, producing the same binary as the original hard-coded
element.

\subsection{Scope of Changes}

\begin{itemize}
  \item \textbf{New file:}
    \texttt{src/elements/BeamKinematicPolicy.hh} ---
    \texttt{beam::FrameData<Dim>},
    \texttt{BeamKinematicPolicyConcept},
    \texttt{SmallRotation} (88 lines),
    \texttt{Corotational} (255 lines).

  \item \textbf{Modified file:}
    \texttt{src/elements/BeamElement.hh} ---
    template signature extended, \texttt{R\_}/\texttt{length\_}
    replaced by \texttt{frame0\_}/\texttt{frame\_n\_}, five methods
    refactored to delegate to the policy.

  \item \textbf{No test files modified:} existing instantiations
    pick up \texttt{SmallRotation} by default.
\end{itemize}

\subsection{Verification}

\begin{itemize}
  \item \textbf{Build:} 44 recompiled targets, zero warnings
        (\texttt{-Werror}).
  \item \textbf{CTest:} 45/47 passed (same 2 pre-existing failures:
        \texttt{cantilever\_benchmark} and
        \texttt{snes\_gmsh\_pipeline}, both missing \texttt{.msh}
        files).
  \item All beam-path tests passed:
        \texttt{beam\_element},
        \texttt{beam\_n\_element},
        \texttt{rc\_fiber\_frame\_nonlinear},
        \texttt{nonlinear\_validation}.

  \item \textbf{Note on \texttt{ShellElement}:}
    The shell element (\texttt{MITC4} Mindlin--Reissner) was not
    refactored in this phase.  A corotational shell formulation
    requires additional machinery (mid-surface rotation
    parametrization, drilling DOF stabilization) and is deferred to
    a future phase.
\end{itemize}
